\subsubsection{Decimal alphabet and check-digit congruence of the emitter}
\label{sec:rec27-codigo}

Retain the emission law \eqref{act21:tpk:bloque-decimal},
with the APP readings and observable \(\ell_9\) of
\eqref{act21:tpk:logaritmo}, the cursors of
Section~\ref{act21:tpk:operaciones}, and the six-window calendar of
Section~\ref{act21:tpk:emision}. The accumulators are \(S_m,E_m,Z_m\),
with \(Z_m=3\), and the output digits are denoted
\(d_{m,1},d_{m,2},d_{m,3}\).
Proposition~\ref{act21:tpk:recuperacion-firmas} already recovers both
residual signatures. Combining this recovery with the actual accumulator
images determines the following alphabet.

\begin{proposition}[Check-digit congruence for the third digit]
Every emitted block satisfies
\[
 d_{m,3}=[5d_{m,2}-2d_{m,1}+1]_{10}.
\]
\end{proposition}
\begin{proof}
The inverse in Proposition~\ref{act21:tpk:recuperacion-firmas}
gives \([E_m]_{10}=[d_{m,2}-d_{m,1}]_{10}\). Substituting
$E_m\equiv d_{m,2}-d_{m,1}\pmod{10}$ and $Z_m=3$ into the third digit gives
$3d_{m,1}+5(d_{m,2}-d_{m,1})+21\equiv5d_{m,2}-2d_{m,1}+1\pmod{10}$.
\end{proof}

\begin{proposition}[Decimal alphabet of forty-two blocks]
The per-window image of the emitter consists exactly of the blocks
\[
\begin{array}{c|rrrrrrr}
[E_m]_{10}&\multicolumn{7}{c}{U_m}\\ \hline
0&114&227&443&556&669&885&998\\
1&129&232&458&561&674&890&903\\
3&149&252&478&581&694&810&923\\
5&169&272&498&501&614&830&943\\
7&189&292&418&521&634&850&963\\
9&109&212&438&541&654&870&983
\end{array}
\]
\end{proposition}
\begin{proof}
A cardinal trajectory reads three consecutive terms of the additive cycle
\[
 (2,3,4,5,6,7,8,9,1).
\]
Their sums are
$(9,12,15,18,21,24,18,12,6)$; reversing direction preserves the set of sums.
Hence
\[
 \operatorname{Im}S_m=\{6,9,12,15,18,21,24\},\qquad
 \operatorname{Im}[S_m]_{10}=\{1,2,4,5,6,8,9\}.
\]
In the multiplicative channel, a fixed factor divisible by three contributes
zero. For invertible factors, the cyclic sums of three readings are
\[
\begin{array}{c|c}
\text{fixed factor}&\text{possible sums}\\\hline
1,8&\{1,3,7,9\}\\
2,7&\{3,5,9\}\\
4,5&\{1,5,7\}
\end{array}
\]
Thus $\operatorname{Im}E_m=\{0,1,3,5,7,9\}$.
Attainability can already be verified in the first window: every cursor in
the following two tables has direction $E$.
\[
\begin{array}{c|rrrrrrr}
S_1&6&9&12&15&18&21&24\\\hline
(i_+,j_+)&(0,8)&(0,0)&(0,1)&(0,2)&(0,3)&(0,4)&(0,5)
\end{array}
\]
\[
\begin{array}{c|rrrrrr}
E_1&0&1&3&5&7&9\\\hline
(i_\times,j_\times)&(2,0)&(0,0)&(0,1)&(1,1)&(0,2)&(0,4)
\end{array}
\]
Independence of the two cursors realizes every one of the $7\cdot6$ pairs.
The first two digits recover the residual pair, so their 42 images are
distinct. The check-digit congruence fixes the third digit and yields
exactly the stated table.
\end{proof}

The decimal block therefore preserves both residual signatures through an
injective encoding and a check digit. This result describes the finite
emitter alphabet; continuation of the regional readings additionally retains
the phase, orientation and memory conditions of the enriched state.

The per-window image therefore has \(42\) blocks. Compatible concatenation
of six windows preserves the census of \(468\) emissions in
Proposition~\ref{act21:tpk:468}, whose ternary reduction yields
\(243\) words within the \(729\)-word ambient space. Each cardinality
counts the object of its stage and retains the provenance of the same
emitter.
